% Quiz 05 --- Modularity in noise, and the likelihood of the degree-corrected
% block model. Printed, worked in class, 20 minutes. Page 1 is every question;
% page 2, printed on the back, is hints for Question 2.
%
% Question 1 asks for a trend and its reason: how the best modularity of a
% random network moves with its average degree. Measured with Leiden on
% Erdos--Renyi networks of N = 10,000 (three draws each):
%   <k>   = 1     2     4     8     16    32    64    128
%   Q_max = 0.99  0.85  0.56  0.36  0.24  0.17  0.11  0.08
%
% Question 2 is the Karrer--Newman (2011) derivation: write the log-likelihood
% in c, theta and w, fix c, find theta-hat and w-hat, substitute them back.
% Page 1 gives the model, the normalization of theta (without it theta and w
% trade scale and the estimates are not unique), the large-N convention that
% removes the self-pairs, and the notation the answer is written in. It gives
% nothing of the route.
%
% Page 2 gives the route as moves, not results. theta-hat = k_i / kappa_r and
% w-hat = m_rs are printed nowhere, because finding them is part (b). Do not
% put them back.
%
% Compile with: xelatex -interaction=nonstopmode quiz05.tex   (run twice)
\documentclass[a4paper, 11pt]{extarticle}
% Shared preamble for Quiz 05 (quiz05.tex and solutions.tex).
% Compile with: xelatex -interaction=nonstopmode <file>.tex   (run twice)
%
% Same house look as the pen-and-paper sheets in
% lecture-note/*/pen-and-paper/: Charter, A4, generous leading, room to draw.

\usepackage[top=0.7in, bottom=0.6in, left=0.8in, right=0.8in]{geometry}
\usepackage{amsmath}
\usepackage{amssymb}
\usepackage{array}
\usepackage{graphicx}
\usepackage[hidelinks]{hyperref}
\usepackage{tcolorbox}
\usepackage{enumitem}
\usepackage{etoolbox}
\usepackage{tikz}
\usepackage{fontspec}
\usetikzlibrary{calc,positioning,arrows.meta}
\definecolor{pltblue}{HTML}{1F77B4}
\definecolor{rulegrey}{HTML}{6A6D75}

% Body face. Charter, and the four ways to get it. XCharter *is* Charter and
% ships with a full TeX Live, so it is asked for first and by file: a CI runner
% answers yes to \IfFontExistsTF{Charter}, then cannot typeset it.
\IfFontExistsTF{XCharter-Roman.otf}{
  \setmainfont{XCharter}[
    Extension      = .otf,
    UprightFont    = *-Roman,
    BoldFont       = *-Bold,
    ItalicFont     = *-Italic,
    BoldItalicFont = *-BoldItalic]
}{
  \IfFontExistsTF{Charter}{\setmainfont{Charter}}{
    \IfFontExistsTF{Iowan Old Style}{\setmainfont{Iowan Old Style}}{
      \IfFontExistsTF{Georgia}{\setmainfont{Georgia}}{}}}}

\setlength{\parindent}{0pt}
\setlength{\parskip}{0.4em}
\setlength{\emergencystretch}{1.2em}
\raggedbottom
\pagestyle{empty}

% ---------------------------------------------------------------------------
% The submission link.
%
% Both the QR square and the printed address are the course's own short link,
% not the form's: paper outlives any one form. Caddy on the course droplet
% points go.skojaku.com/ans-quiz05 at the current Google Form
% (/etc/caddy/Caddyfile, block go.skojaku.com). Swapping the form is a line on
% the server, not a reprint.
%
% If the link ever changes, rebuild the square with
%
%   uvx --from segno segno --output=quiz05-form-qr.png --scale=20 --border=1 \
%       --error=m "<the url>"
%
% PNG rather than PDF for the square: segno's PDF is too spare for xdvipdfmx to
% include, and a QR is squares, so 700px across a 3cm box is already sharper
% than the printer.
% ---------------------------------------------------------------------------
\newcommand{\formlink}{https://go.skojaku.com/ans-quiz05}
\newcommand{\formshort}{\href{\formlink}{\texttt{go.skojaku.com/ans-quiz05}}}

% Quiz header. #1 = title, #2 = the line under it, #3 = term.
\newcommand{\quizhead}[3]{%
  {\footnotesize\textsc{Advanced Network Science} \hfill \textsc{#3}}\par
  \vspace{-0.2em}\rule{\textwidth}{1pt}\par
  \vspace{0.3em}
  {\LARGE\bfseries #1}\par
  \vspace{0.1em}
  {\small #2}\par
  \vspace{0.7em}
  {\small Name \underline{\hspace{6cm}} \hfill Date \underline{\hspace{3.2cm}}}\par
  \vspace{0.7em}}

% A question heading with its point value.
\newcommand{\question}[3]{%
  \vspace{0.5em}
  {\large\bfseries Question #1 \quad\normalsize\mdseries (#2 points)}\par
  \vspace{-0.3em}{\color{rulegrey}\rule{\textwidth}{0.6pt}}\par
  \vspace{0.1em}
  {\bfseries #3}\par
  \vspace{0.2em}}

% An empty box to draw or write in. #1 = height, #2 = caption under the rule.
\newcommand{\drawbox}[2]{%
  \begin{tcolorbox}[colback=white, colframe=rulegrey, boxrule=0.5pt,
                    sharp corners, height=#1, valign=top,
                    boxsep=2pt, left=4pt, right=4pt, top=3pt, bottom=3pt]
  {\footnotesize\color{rulegrey} #2}
  \end{tcolorbox}}

% Ruled writing lines, #1 of them.
\newcommand{\writelines}[1]{%
  \vspace{0.3em}
  \foreach \i in {1,...,#1}{%
    {\color{rulegrey}\rule{\textwidth}{0.4pt}}\par\vspace{1.5em}}}


% As in Quiz 4: title and points share one line, to keep page 1 to one page.
\renewcommand{\question}[3]{%
  \vspace{0.4em}
  {\large\bfseries Question #1 \quad #3\quad\normalsize\mdseries (#2 points)}\par
  \vspace{-0.4em}{\color{rulegrey}\rule{\textwidth}{0.6pt}}\par}

\begin{document}

\quizhead{Quiz 5 \; --- \; Communities: noise and likelihood}
        {20 minutes \quad$\cdot$\quad 10 points \quad$\cdot$\quad
         Closed notes, work on your own \quad$\cdot$\quad
         hints for Question 2 on the back}
        {Fall 2026}

\begin{tcolorbox}[colback=white, colframe=rulegrey, boxrule=0.5pt,
                  sharp corners, boxsep=1pt, top=4pt, bottom=4pt]
\begin{minipage}[c]{0.14\textwidth}
\begin{center}
\includegraphics[width=1.6cm]{quiz05-form-qr.png}
\end{center}
\end{minipage}
\hfill
\begin{minipage}[c]{0.83\textwidth}
{\small
You hand in two photos, one per question. Write each question's answer on
its own page, photograph each page, and upload it. Point a camera at the
square, or type \formshort.
}
\end{minipage}
\end{tcolorbox}

% ---------------------------------------------------------------------------
\question{1}{4}{Communities in pure noise.}

You generate Erd\H{o}s--R\'enyi random networks with $N = 10{,}000$ nodes:
every pair of nodes is joined with the same probability, so no communities are
planted. You make one at each average degree
$\langle k \rangle = 1, 2, 4, \dots, 128$, and on each you find the partition
with the highest modularity, $Q_{\max}$.

(a) [2 pts] As $\langle k \rangle$ grows, does $Q_{\max}$ go up, go down, or
stay flat? Sketch $Q_{\max}$ against $\langle k \rangle$ (log scale for
$\langle k \rangle$). What value does $Q_{\max}$ head toward when the network
is very sparse, and when it is complete ($\langle k \rangle = N - 1$)?

(b) [2 pts] Explain why $Q_{\max}$ moves this way. Then: a paper reports
$Q = 0.45$ for a network with $N = 10{,}000$ and $\langle k \rangle = 4$, and
calls it strong community structure. What would you check before believing
it?

% ---------------------------------------------------------------------------
\question{2}{6}{The likelihood of the degree-corrected block model.}

\begin{tcolorbox}[colback=pltblue!4, colframe=pltblue!35, boxrule=0.5pt,
                  sharp corners, boxsep=2pt, left=6pt, right=6pt,
                  top=4pt, bottom=4pt]
{\small
An undirected network has $N$ nodes and $M$ edges. $A_{ij} = A_{ji}$ is the
number of edges between nodes $i$ and $j$ ($A_{ii} = 0$), and
$k_i = \sum_j A_{ij}$ is the degree of node $i$. Node $i$ belongs to group
$c_i \in \{1, \dots, K\}$. The \emph{degree-corrected block model} draws, for
every pair $i < j$ independently,
\[
  A_{ij} \sim \mathrm{Poisson}\!\left(\theta_i \theta_j\, w_{c_i c_j}\right),
  \qquad
  \text{with} \quad \sum_{i:\,c_i = r} \theta_i = 1 \ \text{ for every group } r,
\]
where $\theta_i > 0$ and $w_{rs} = w_{sr} > 0$. Write $m_{rs}$ for the number
of edge ends in group $r$ whose other end is in group $s$ (an edge inside
group $r$ counts twice in $m_{rr}$), and
$\kappa_r = \sum_{i:\,c_i = r} k_i$ for the total degree of group $r$.

$N$ is large: you may replace $\sum_{i<j}$ by $\tfrac{1}{2}\sum_i \sum_j$, and
drop any term that depends on none of $c$, $\theta$, $w$.
}
\end{tcolorbox}

(a) [2 pts] Write the log-likelihood
$\mathcal{L}(c, \theta, w) = \log P(A \mid c, \theta, w)$. Simplify it until no
sum over pairs of nodes is left.

(b) [2 pts] Fix the assignment $c$. Find the maximum-likelihood estimates
$\hat{\theta}_i$ and $\hat{w}_{rs}$.

(c) [2 pts] Substitute them back, and write $\mathcal{L}$ as a function of the
assignment $c$ alone, in terms of $m_{rs}$ and $\kappa_r$.

\vfill

{\footnotesize\color{rulegrey}
Done early? (not graded) Without the $\theta$'s, the same steps give
$\tfrac12 \sum_{r,s} m_{rs} \log\!\big(m_{rs} / n_r n_s\big)$, where $n_r$ is
the number of nodes in group $r$. What did $\theta$ change, and why does that
stop the model from putting all the hubs in one group?}

% ===========================================================================
\newpage

{\large\bfseries Hints for Question 2}\par
\vspace{-0.4em}{\color{rulegrey}\rule{\textwidth}{0.6pt}}\par
\vspace{0.2em}
{\small Read only as far as you need. Each hint names a move; working it out is
still yours.}

\vspace{0.6em}

\textbf{Hint 1 --- The Poisson distribution.}
If $A \sim \mathrm{Poisson}(\lambda)$, then
\[
  P(A = a) = \frac{\lambda^{a} e^{-\lambda}}{a!} .
\]
The pairs are independent, so the likelihood is a product over pairs, and its
logarithm is a sum over pairs.

\vspace{0.6em}

\textbf{Hint 2 --- From pairs to nodes.}
A factor that depends on only one end of the pair can leave the inner sum:
\[
  \sum_i \sum_j A_{ij}\, f(i) \;=\; \sum_i f(i) \Big( \sum_j A_{ij} \Big) .
\]
What is the bracket? Since $A_{ij} = A_{ji}$, a factor $f(j)$ goes the same
way.

\vspace{0.6em}

\textbf{Hint 3 --- From pairs to groups.}
A factor that depends only on the groups of the two ends can be collected one
pair of groups at a time:
\[
  \sum_i \sum_j A_{ij}\, g(c_i, c_j)
  \;=\; \sum_r \sum_s g(r, s) \sum_{i:\,c_i = r} \ \sum_{j:\,c_j = s} A_{ij} .
\]
What does the last double sum count? Make the same move on the term that has
no $A_{ij}$ in it, and then use the normalization of $\theta$.

\vspace{0.6em}

\textbf{Hint 4 --- Maximizing under a constraint.}
The $\theta_i$ of one group must sum to 1, so you cannot just set
$\partial \mathcal{L} / \partial \theta_i = 0$. Add one Lagrange multiplier
per group and maximize
\[
  \mathcal{L} \;-\; \sum_r \mu_r \Big( \sum_{i:\,c_i = r} \theta_i - 1 \Big) .
\]

\end{document}
